% E4 (SCRUM-25): Scope and Limitations.
% Scope = the core in lepton-plan.md's overview, as built: decision-tree row "QC weekly counts
% by onset week", option C, g = 2 (Notes.md 1.5); Kristine added as a fourth storm (Notes.md
% 1.1; confirmed for E4 on 2026-10-21). Limitations = every item the plan cut, plan corrections
% 1, 2, 8 and 9, the assumptions table (tables/assumptions.tex), Notes.md sections 1-3 and the
% decision tree's "say in the paper" line for option C.
% Numbers: results/forecasts.csv (454 forecasts per model, 349 in the MAE window, M3's 1,065
% for 2018 week 31, the 1.5-case threshold), results/c8_observed_runs.csv (Carina 120.0 mm,
% R0 119.7 mm, 1,414 exposed), results/m3_reproduction.md, data/raw/qcesd.md (207 published,
% 264 in the list, 17 days), data/raw/noah.md (11 barangays, 542 people), Notes.md 1.5
% (83, 163, 176 and 775 cases), results/m4_residual_acf.csv (M4's misses 1 and 2 weeks apart:
% 0.41 and 0.15, 95% band 0.08 and 0.09; python -m src.m4), config.py.
% Scope's first two paragraphs (October 23, at the team's request): LEPTON is designed for any
% Philippine city; QC came first as the most populous city (background.tex, psa_census2024) with
% an up-to-date case list (data/raw/qcesd.md: onset 2015-01-30 to 2026-10-14, received 20 October
% 2026), and because the schedule left time to request counts from one city only. The national
% fallback is the decision tree's last row (lepton-plan.md). NOAH's maps cover most provinces:
% 66 to 79, by return period (data/raw/noah_national.md). Every other source is in Appendix E.
% Future work, M2 and the two-week forecast (October 23): Implementation Plan v2 (9 October
% 2026), p. 16, defines M2 as "Prophet, yearly seasonality only, no regressors"; its pp. 28-29
% ("Observed versus forecast weather") score the two-week forecast on observed weather for the
% coming week, an optimistic upper bound (the plan's "ceiling run"), and on archived forecasts,
% and note that the one-week forecast uses observed weather only. The flag rule (25%, 50%) was
% confirmed by the team on October 23.
\section{Scope and Limitations}\label{sec:scope}

\subsection{Scope}\label{sec:scope-scope}

LEPTON is designed for any Philippine city, since every input but the case counts comes from a
national or global source (Appendix~\ref{app:sources}): IMERG, ERA5 and ECMWF rain, the WorldPop
grid, the PSA census and birth tables, the national barangay boundaries, the DOH facility
registry, the HydroBASINS river basins and Project NOAH's flood maps, which cover most provinces.
Another city therefore needs its own weekly case counts for Layer~1, and its catchments,
barangays and health centers, drawn from the same sources, for both layers; Layer~2 needs no
case data.

This study covers one city, Quezon City, for two reasons. First, QC is the country's most
populous city, and QCESD keeps an up-to-date list of the city's leptospirosis cases with each
case's date of onset, from January 2015 to mid-October 2026, which supports weekly forecasting.
Second, the competition's schedule left time to request case counts from one city only,
and a reply in time was not assured: without one, Layer~1 would have fallen back to national
weekly counts \parencite{parreno2023}, without flood features. QCESD's list arrived on 20 October
2026. Because only one city was tested, results may not carry over to cities with different
drainage, climate or reporting. Cases are counted by week of onset (CDC/MMWR weeks), and no
case-level record is published.

\emph{Layer 1} forecasts next week's number of cases in QC by week of onset, with a 90\% range
and the probability that the count exceeds the epidemic threshold, set here at the same week's
mean plus two standard deviations over up to five recent non-COVID years. Four models are
compared: two seasonal baselines, M0 (the count 52 weeks earlier) and M1 (the same week's
average over the threshold's years); the published Prophet--XGBoost hybrid, M3
\parencite{latigay2026}, first reproduced on its own national data \parencite{parreno2023}; and
the proposed negative binomial regression, M4, also refitted without its flood exposure index to
test whether flood exposure adds accuracy. In a weekly
rolling-origin backtest every model made 454 one-week-ahead forecasts, for the weeks from
14--20 January 2018 to 20--26 September 2026. The primary endpoint is M4's mean absolute error
against M3's over the 349 forecasts outside the COVID-19 years 2020--2021, with a 95\%
moving-block bootstrap interval \parencite{kunsch1989}.

\emph{Layer 2} counts, for each of QC's 142 barangays, the residents in each hazard level of
Project NOAH's 5-, 25- and 100-year flood maps \parencite{noah_nd}, from the WorldPop 2024
population grid \parencite{bondarenko2025} fitted to the 2024 census
\parencite{psa_census2024}. A storm's rain over the river catchments that drain into the city,
set against that rain's return levels, decides how many of those residents it exposes. The
exposed residents then become stock under DM~\mbox{2023-0309} \parencite{doh2023}: doxycycline 100~mg
capsules for adults by risk tier, courses for children and amoxicillin 500~mg tablets for
pregnant women, per barangay and per health center. These are stock-planning figures, not
prescriptions.

Layer 2 was run on four storms: Typhoon Carina (July 2024), Severe Tropical Storm Kristine
(October 2024), the July 2025 floods and the August 2026 habagat floods. Each ran on observed
satellite rain from NASA's IMERG \parencite{huffman2024} and on ECMWF forecasts issued one, two
and three days before its rain, from Open-Meteo's archive \parencite{zippenfenig2023}. Observed
rain is set against IMERG's own return levels and forecast rain against ERA5's
\parencite{hersbach2020}, because the two products differ in scale. Kristine was not in the
project proposal; it was added because IMERG puts Carina's rain at the no-flood floor
(Section~\ref{sec:limitations}), so Carina alone could not show the chain at work. Each run is
compared with the barangays that city and news reports named as flooded. LEPTON is meant as a
decision-support tool for health offices, not a replacement for clinical judgment.

\subsection{Limitations}\label{sec:limitations}

The limitations are grouped by the part of the study they affect. Table~\ref{tab:assumptions},
in Appendix~\ref{app:assumptions}, lists every assumption behind them, with its value and how it
was tested.

\paragraph{Case data}
QCESD's list gives each case's dates of onset and admission but not the date the case was
reported, so what the office knew on each forecast day cannot be rebuilt. Cases are therefore
forecast by week of onset from final counts and, at each forecast, the last two weeks of counts
are hidden from every model, as if they had not been reported yet. Two weeks is an assumption,
since QCESD's reporting delay is not known. The backtest therefore uses final counts with a
two-week guard, and real-time accuracy may be lower. Counts published during a surge show the
gap: QCESD's total for 1 January to 3 August 2025 was 207 cases \parencite{herriman2025}, while
the final list holds 264 cases admitted by then. A week counts as final once 17 days have
passed since it ended, because 99\% of cases were admitted within 17 days of onset; the delay
from admission to report is unknown, so this is a lower bound, and the weeks after 20--26
September 2026 are left out as still filling in. Without report dates, those weeks also cannot
be nowcast, a step a live system would need.

Surveillance data likely undercount leptospirosis. QC's early years are also small: 83, 163 and
176 cases in 2015--2017, against 775 in 2018. If fewer facilities reported then, the thresholds
and M1 forecasts that average over those years are too low; the first threshold of 2018 is 1.5
cases a week. The list's own morbidity-week column changes meaning between years, so every
week is taken from the date of onset. The COVID-19 years 2020--2021 carry an indicator in M4, are
not threshold reference years and are left out of headline scores; all years are also scored,
as a check (Section~\ref{sec:results-checks}). Leaving them out joins 2019 directly to 2022 in the block bootstrap, once in 349
weeks. Finally, large surges are few, and the three largest came in the last three years, so no
model has many outbreaks to learn from.

\paragraph{Rain and flood exposure}
Satellite rain under-reads typhoons. In five typhoon rain events in the Philippines, IMERG
considerably underestimated the rain that gauges recorded \parencite{aryastana2022}, and
calibrating it against Philippine gauges brings it closer to them \parencite{veloria2021}.
Carina shows the cost: on
IMERG, its wettest three days over the catchments total 120.0~mm, at the no-flood floor $R_0$
of 119.7~mm, so its observed-rain run exposes about 1,400 residents and sizes almost no stock,
even though QC's cases surged one to two weeks later. Gauge-calibrated products track
Philippine typhoon rain more closely \parencite{bagtasa2022}, but switching would mean a new
26-year download, and IMERG's Final run arrives about three and a half months late, too late
for August 2026. ERA5, whose return levels judge the forecasts, also underestimates the highest
observed totals, with its largest errors in the tropics \parencite{lavers2022}; because each product is
judged on its own scale, this moves where a forecast falls but favors neither run.

A storm is measured by its wettest three consecutive days of IMERG rain over the catchments,
from the daily product on UTC days rather than a rolling 24-hour maximum, which would need
about 440,000 half-hourly files. A single wettest day missed the long habagat floods of 2025
and 2026, so three-day totals are used and the single day is kept as a comparison. NOAH's maps do
not state the duration of the rain they were drawn for, so neither choice can be tied to the
maps. Return levels rest on 18 to 26 years of IMERG, depending on the year: in a simulation with
25 years of annual maxima, 90\% of fitted 25-year levels fell between 16\% below and 14\% above
the true value, which moves where a storm falls between the maps, and fewer years widen that
range. Below the floor $R_0 = 0.5\,\hat{R}(5)$, half the five-year rain, nobody counts as
exposed; it is a starting value, left untuned because tuning needs floods that the city's
disaster office has listed.

The flood maps are static scenarios, not observed floods. NOAH's 100-year map is a later release
whose flood shapes are thinner in places: in 11 barangays a rarer map puts fewer residents in
flood zones than the five-year map does, by at most 542 people (San Jose, on the 100-year map).
The maps are used as published. WorldPop's counts are modeled, from a release (R2025A) that
WorldPop calls an alpha, and the 2020 census's age shares are applied to 2024 totals. The flood
exposure index is therefore a proxy for exposure, not a measure of it.

\paragraph{Models and evaluation}
The benchmark's published gain could not be reproduced. Rebuilt on the benchmark's own national
data, the seasonal naive forecast matched all eight of the paper's published scores for it, so
the data and the scored weeks agree, and Prophet came within 4.4\% of its published mean
absolute error (24.73 against 23.68). The rebuilt hybrid's error, however, was 24.60 against
the published 16.72, about the same as Prophet alone, and none of the 75 XGBoost settings tried
came within 10\%. Training XGBoost also on the forecast weeks' own errors, which a real forecast
cannot know, gives errors like the paper's; whether that happened cannot be told without the
authors' code. M4 is therefore compared with the reproduction of M3, and 16.72 is not quoted as
M3's accuracy. Ported to QC, M3 is also fragile early in the backtest, when only 2015--2017 lie
behind it: it forecast 1,065 cases for the week of 29 July--4 August 2018, which had 58.

The epidemic threshold is a choice made for this study, and QCESD's method may differ. Only
one-week-ahead forecasts are made. M4's misses come in runs: in its fit on the last forecast day,
one week's miss is correlated with the next week's (0.41) and with the one after (0.15), both
outside the 95\% range expected if weeks were independent ($\pm0.08$ and $\pm0.09$), consistent
with the newest count M4 can read being three to four weeks older than the week it forecasts; the
moving-block bootstrap allows for such runs. Interval and breach scores start in 2022, because each model's forecast spread is
estimated from its own past forecasts, which need a warm-up.

\paragraph{Prophylaxis stock}
Stock is sized by the regimens of DM~\mbox{2023-0309} \parencite{doh2023}. Its moderate-risk
group, adults with a wound or with floodwater in the eyes, nose or mouth, cannot be counted from
any available data, so for short floods the share $w$ of exposed adults in it is reported at 0,
0.5 and 1; the first and last are bounds. A flood lasting two days or more makes every exposed
adult high risk, this study's reading of continuous exposure, tested at one and three days. The memo makes doxycycline, 4~mg
per kg up to 200~mg, the first-line regimen for every child, so children are counted in courses
that a pharmacy turns into tablets by weight. Pregnant women's stock is amoxicillin 500~mg twice
a day for three days, from section II.B.2.b of the memo; PhilHealth's 2012 recommendations give
no prophylaxis regimen for pregnant women \parencite{philhealth2012}, so this regimen rests on
the memo alone, and women allergic to penicillin, who take erythromycin instead, are not stocked
separately. Lactating women are not counted apart from other adults: the census tables cannot
count them, and the memo gives them no regimen.

Pregnant women are estimated as 0.75 times the registered births to mothers living in QC, split
among barangays by their share of women aged 15--49; the registrations are not adjusted for
under-registration \parencite{psa_births2026}. Every resident of a flooded zone is counted as
touching the water, at any depth, and stock is reported at full coverage; a health office scales
the totals by the share of residents it expects to reach. Flood durations come from city and news
reports and may be off by a day (tested at $\pm1$ day); Kristine's one-day flood rests on
reports from 24 October only. A forecast cannot know how long the flood will last, so the
forecast runs report a one-day and a seven-day flood. Nor can it know when the rain will peak:
each lead takes the forecast's wettest three-day window among those ending on the storm's
observed wettest days, a choice made with hindsight. A live system would scan every window the
forecast covers and could pick other days and another amount, so the forecast runs show what
the forecasts said about the storm's own days, not the trigger an office would have seen.
Stock goes to the 66 health centers that
the national facility registry places inside their own barangays; the City Health Office is
treated as holding the city's stock rather than handing it out.

Finally, Layer 2 is checked against flood reports, not against cases. The reports name 8 to 17
flooded barangays per storm and are incomplete, and a barangay counts as flagged when at least
25\% of its residents live in zones the rain floods (50\% is also reported), a line chosen
because every barangay has some residents in the five-year zones. The storm runs are therefore
demonstrations, not
validations: no barangay case counts were used, and four storms are too few to measure forecast
error. LEPTON estimates stock needs, not supply chains, delivery or uptake, and every dose
figure in this paper is a stock-planning figure, not a prescription: DM~\mbox{2023-0309} requires a
physician's consultation before any prophylaxis.

\subsection{Not Covered: Future Work}\label{sec:future}

To finish within the competition's schedule, several parts of the project's original plan were
cut. They are named here so that no result is read as covering them:
\begin{itemize}
  \item other Philippine cities, Cebu City first, each with its own weekly case counts
    (Section~\ref{sec:scope-scope});
  \item running LEPTON live each week, which needs counts as they arrive and report dates to
    nowcast the weeks still filling in;
  \item forecasting cases two weeks ahead. The one-week forecast reads only rain that has
    already fallen; a two-week forecast could also use the coming week's rain from a rain
    forecast, and would then be scored twice, on the rain that fell (an optimistic upper bound)
    and on the archived forecasts that existed on each forecast day;
  \item two more models for comparison: Prophet with yearly seasonality only (M2) and an
    XGBoost model (M5);
  \item other population and flood maps: Meta's population map and the LiDAR-based flood maps
    on LiPAD;
  \item the full sensitivity analysis, of which only the tests in Table~\ref{tab:assumptions}
    were run;
  \item tuning the no-flood floor $R_0$ on floods the city's disaster office has listed;
  \item a distributed-lag non-linear model of rain and cases;
  \item checking flagged barangays against satellite flood maps (Sentinel-1) and against
    barangay case counts; and
  \item forecasting hospital demand.
\end{itemize}
